\documentclass[sigconf, 10pt]{acmart}
\begin{document}

\title{Coded Bias Ethics pooppooppooppooppooppooppooppooppooppooppooppooppoop}
\author{Roman Marquez-Twisdale}
\email{r.marqueztwisdale@digipen.edu}

\date{Today}
\maketitle

\section{Background Information}
In this paper, we will review the film Coded Bias, dive deep into the technical subjects being focused on, and analyze the ethics of how such technologies are being used in the world today, particularly those focused on by the film. “Coded Bias follows M.I.T. Media Lab computer scientist Joy Buolamwini, along with data scientists, mathematicians, and watchdog groups from all over the world, as they fight to expose the discrimination within algorithms now prevalent across all spheres of daily life.” ~\cite{used_twelfth(11)} Specifically, we will discuss the discriminatory effects of such tools in the realms of law enforcement, employment, and financial opportunity. The following actors represent the relevant powers which make such systems of discrimination possible.
\subsection{The Public}
Commonwealth should not be underappreciated in any context. Concepts like “society”, “government”, “nation”, and many others within the realm of civil matters are intimately dependent on the existence of a general population. The technologies being discussed in Coded Bias are tools to be used by a society, government, or nation which are systems implemented by the public. And the victim of the harm exposed and highlighted by the documentary is once again, the public.
\subsubsection{Role of the Public}
The public is the medium upon which the ML tools being discussed in Coded Bias are used. It is the role of the public to submit to the collection of their data for such tools to be trained and used, or to push back against their use. It is the job of the people to pay attention to the developments taking place in their home, and to react accordingly. It is also the people who should benefit from the addition of new technologies in their societies. Ideally, there should be system for measuring the success of new technologies. Without the ability to quickly analyze the effects of a new methodology, how can the people know whether it is working? The implementation of machine learning algorithms and models in areas such as law enforcement, employment, and financial risk assessment has not yet been explored. Measuring the impact such models are having in these sections of society will most likely take time.
\subsection{Law Enforcement Organizations}
One of the most potent forms of discrimination discussed in the documentary stems from the use of ML in law enforcement. Policing and law enforcement agencies/organizations are supposed to uphold the law. In a country where the people have pushed for legislation that supports equal opportunities for all, this includes protection against discrimination.
\subsubsection{Role of Law Enforcement Organizations}
It’s worth mentioning that most law enforcement organizations are directly supported by the public they serve. For example, in the United States, “In 2021, local governments provided approximately 87\% of total police department funding.” ~\cite{used_eighth(7)} They are truly a service for the people, by the people. To enforce any law, there must first be a clear description of said law, and the enforcing agent must have a clear understanding of the legislation they are attempting to enforce. Machine learning models and tools have been advertised as increasing efficiency and broadening capabilities *in general*. It is one of their great strengths. It makes sense that a law enforcement agency might seek their abilities to better uphold the law. This would be well within the realm of appropriate interests for a law enforcement agency to have. 
\subsection{Tech Industry}
Large tech companies responsible for developing and selling AI models across all industries make the exposed harm possible. These organizations control vast amounts of wealth and influence in the world economy. As of February 2026, roughly 32.41\% of the S\& P 500 index represented the information technology industry. With such profound influence in the global economy, these organizations are under a lot of pressure to maintain a financially stable existence.
\subsubsection{Role of Tech Industry}
What should we expect from these large tech companies? According to the Friedman essay, a Harvard law document published in the New York Times, “The purpose of a corporation is to conduct a lawful, ethical, profitable and sustainable business in order to ensure its success and grow its value over the long term.” ~\cite{used_ninth(8)} It is quite clear that these businesses are pursuing financial growth and stability, but what about the lawful and ethical portions of this description? “This requires consideration of all the stakeholders that are critical to its success (shareholders, employees, customers, suppliers and communities), as determined by the corporation and its board of directors using their business judgment and with regular engagement with shareholders, who are essential partners in supporting the corporation’s pursuit of its purpose.” ~\cite{used_ninth(8)} Being considerate of *communities* is likely to raise questions about the potential harm that could be caused due to the hasty implementation of new technologies in a world without sufficient regulation for those technologies.
\subsection{Legislators}
Legislation defines how to operate in a society. Whether that’s a business, a single citizen, or a country. Legislators around the world have the power to support the people they represent and stand for. Without legislation, the road to deciding ethical conduct for any entity is long and messy.
\subsubsection{Role of Legislators}
Legislators in the field of technological development and implementation should be knowledgeable about the industries they write legislation for. This means staying up to date with the capabilities of machine learning algorithms, their strengths, and their weaknesses. This means putting thought into developing legislation which will allow the technology to develop and grow, whilst ensuring public safety.
\section{The Coded Bias Problem}
Coded Bias explores the weaknesses of machine learning algorithms and the harm being caused by the hasty implementation of such algorithms in various sectors of society. But ultimately focuses on three main observations.
\subsection{Algorithmic Bias and Discrimination}
Algorithmic bias exists and is not just possible, but likely to form unless intentionally avoided. Even when actively seeking bias and attempting to avoid it, removing all bias from a machine learning system might not be possible. To be responsible with this technology, we must recognize this risk, and act accordingly. Bias in these models has many origins.
\subsubsection{Training Data}
“The data used to train an LLM may be drawn from a non-representative sample of the population, which can cause the model to fail to generalize well to some social groups. The data may omit important contexts, and proxies used as labels (e.g., sentiment) may incorrectly measure the actual outcome of interest (e.g., representational harms).” ~\cite{used_first(0)} The collection of data might also introduce bias. Unless the data is already well understood, it’s possible to manifest an overgeneralization of distinct social groups. Majorities within the data could overshadow less frequent details. “Of course, even properly collected data still reflects historical and structural biases in the world.” ~\cite{used_first(0)}
\subsubsection{ML Algorithms}
We must also consider our use of such messy data. “The training or inference procedure itself may amplify bias, beyond what is present in the training data. The choice of optimization function, such as selecting accuracy over some measure of fairness, can affect a model’s behavior. The treatment of each training instance or social group matters too, such as weighing all instances equally during training instead of utilizing a cost-sensitive approach.”~\cite{used_first(0)} Once we have trained a model, we must consider how we test that model as well. A perfectly balanced model (if such a thing could even exist) could still be warped by our opinions and resulting decisions about how to implement that model or what to use it for.
\subsubsection{Testing Methods and Standards}
Aside from our own biases, the standards which guide our testing of models is subject to bias as well. Datasets used to measure success might fail to accurately represent the population which the model will be used on. “The choice of metric can also convey different properties of the model, such as with aggregate measures that obscure disparate performance between social groups, or the selection of which measure to report (e.g., false positives versus false negatives).” ~\cite{used_first(0)}
\subsubsection{Deployment}
On top of all the ways in which a model might be created with bias, there is also the question of how we choose to use it. We may choose to implement a model within a setting that is drastically or even slightly different from the environment it was intended to operate in. We may implement a model behind an interface which introduces its own bias, which is effectively the same as a faulty environment. There are many ways to introduce bias. When it comes to automated systems, any existing bias becomes much more potent and inflammatory.
\subsection{Use of Facial Recognition in Law Enforcement}
In the film, the main example used to represent the harms of bias automation is that of law enforcement in the United States. “Seven law enforcement agencies within the Departments of Justice (DOJ) and Homeland Security (DHS), such as the Federal Bureau of Investigation and U.S. Secret Service, reported using facial recognition technology to support criminal investigations.” ~\cite{used_second(1)} It was similarly reported that all seven of these organizations employed systems created by other entities including some that were privately owned. Legal systems are delicate and complicated. The addition of such untested tools in such inflammatory environments is a recipe for disaster. The documentary highlights the dangers to people’s right to privacy. As we will discuss in the next section, the suggested solution to this risk is to tighten legislation around the use of automated systems. 
\subsection{Lack of Regulation for Machine Learning and Algorithmic Recognition Software}
In the United States, the Government Accountability Office found that three of the seven law enforcement organizations discussed in the previous section had policies or standards in place for how to respect civil rights and liberties while employing facial recognition systems. But the others, and many more clients of automation do not have such guidelines for using the new tools. Considering the issue at hand is legislation, it’s even more alarming that organizations like the DOJ don’t rigorously imposed policies on how to use automation. Coded Bias urges the world to see this and testifies before congress to push for legislation for facial recognition in the United States.
\section{Technical Background}
Coded Bias highlights the problems with careless training and implementation of automation in general, but the focus of the film is on the effects of machine learning models trained for facial recognition and automated decision-making systems. In this section we will pick apart said models and define the specific fields where automation can cause obvious harm.
\subsection{CNNs for Facial Recognition}
Machine learning is a broad field of computer science. The specific branch of ML which allows for sophisticated surveillance techniques such as facial recognition is called deep learning. Deep learning utilizes massive structures of artificial neurons known as neural networks. There are two main types of neural networks used today. Fully Connected Neural Networks and Convolutional Neural Networks(CNN). CNNs are known for their enhanced ability to process and recognize visual patterns within images and videos.
\subsubsection{Face Detection}
Object detection represents the ability to detect the sheer presence of a specific object within an image or video. This ability lays the foundation for more sophisticated functionality. This usually means finding the specified object within the image or video frame and defining a bounding box which encapsulates the entire entity. Face detection is simply using object detection to find faces. “It is the first and probably the most essential step for facial recognition, and is used to extract human faces from digital images.” ~\cite{used_eleventh(10)}
\subsubsection{Face Analysis}
Face analysis is the process of extracting meaning from a detected face. The goal is to obtain other kinds of information belonging to the person behind the face. “…face analysis doesn’t pursue identification or verification. Instead, it focuses on gathering actionable insights from facial expressions and positioning, without putting a name to the face.” ~\cite{used_third(2)} Image segmentation datasets are used extensively when training face analysis models. “Image segmentation involves labeling each pixel of an image with a class.” ~\cite{used_fourth(3)} The ability to classify each pixel is essential when attempting to extract details from images or video frames.
\subsubsection{Facial Recognition/Classification}
Face recognition models are another way to use face detection. “Where face detection simply identifies the presence of a face in an image, face recognition either verifies or identifies an actual person.” ~\cite{used_third(2)} This is one of the most essential capabilities for a model to have if it is going to be useful in surveillance. CNNs utilize special kinds of weights known as filters. Each filter is a small matrix of weights which make passes over the input data (usually an image or a video frame). During each pass it calculates a sum of its weights and the RGB values associated with the pixels it is currently covering. This sum represents a single output. After the matrix has passed over the entire image input, another image-like output has been generated which will be fed into the next convolutional layer. After a few of these layers process the original image input, the output is then flattened and fed into a traditional fully connected neural network.
\subsubsection{Training Datasets for Computer Vision}
Training datasets for these models consist of images. The images used for training the models discussed in Coded Bias would have been images of people or people’s faces.
\subsection{Automated Decision-Making Systems}
Automated decision-making systems represent one of the most direct ways in which bias can cause harm. Many areas of human life have become saturated with automation due to its obvious benefits, but Coded Bias focuses on three main examples.
\subsubsection{In Policing}
Law enforcement and Policing take the main spotlight within the film. Within this area alone, there are many facets in which automation can be implemented. Automation has many uses in identification and surveillance “From recognizing faces, fingerprints, and other biometric identifiers, to tracking license plates and locating gunshots, AI has a wide range of existing and potential applications for identification and surveillance in criminal justice contexts.” ~\cite{used_fifth(4)} Of course, as is the focus of the film, when the use of such tools goes awry, the damage can be serious. False identification can lead to mistaken arrests. Forensic analysis is benefitting from machine learning as well. The speed and accuracy of certain models can greatly benefit the process of analyzing crime scene data. But the potential to reach false conclusions in such a delicate and impactful system is potentially as much of a risk as the crimes being dissected. Predictive policing also makes use of automation. “Law enforcement agencies use historical data to forecast the places where crime is likely to cluster and people who are at a higher risk of engaging in or being victims of criminal activity.” ~\cite{used_fifth(4)} The fairness of automated system behavior can be compromised by the quality of the training data. If these algorithms learn from data representing a history of disproportionate policing, the model will amplify those biases, continuing these harmful imbalances. The last area we will discuss is risk assessment. Tools for risk assessment estimate the chance that specific outcomes will occur in the criminal justice system. These tools are widely used to inform pretrial release, sentencing, prison classification, probation, parole, and supervision. ~\cite{used_fifth(4)} Obviously this can have negative side effects for anyone who is wrongly assessed.
\subsubsection{In Hiring and Employment}
Employment is another huge area adopting the use of machine learning models. Automated systems in this area are used for a wide range of decision-making processes such as recruitment, hiring, and deciding merit for promotions. While these tools can bring myriad benefits, they can also exacerbate existing biases and contribute to discriminatory outcomes. ~\cite{used_sixth(5)} Minorities used to get the short end of the stick may be harmed even more by such systems. Humans are not perfect, but at least we cycle in and out of the workplace. An automated decision-making system could potentially be implemented and then left alone if it seems to be working for an unknown and potential extremely long period of time before anyone notices that it was favoring a specific demographic. The bias in such a system may never be noticed due to the subtlety of its effects. This is arguably one of the most dangerous outcomes.
\subsubsection{In Financial opportunity and Credit}
A machine learning model used in a financial credibility system would reduce time spent on many tasks. “It automates document verification, income validation, fraud detection, and risk scoring in a single workflow, which improves operational efficiency.” ~\cite{used_seventh(6)} These models will be used for predictive analysis tasks such as diving into a subject’s historical loan performance and borrower data and then using that information to potentially change interest rates or loan accessibility. “It can also evaluate transaction behavior, spending habits, and repayment trends.” ~\cite{used_seventh(6)}
\section{Ethical Analysis}
This section will analyze the ethical implications of three main questions regarding the focus points of Coded Bias. The analysis be through the lens of social contract theory. I believe this ethical framework is most appropriate due to one of the most prominent sources of motivation for adoption of such systems of automation; public benefit.
\subsection{Question One}
The first of the three questions is “Is automation ethical in the context of threat detection?” Social contract theory assumes no individual is above the law. This is part of what keeps laws just, for why would anyone write an unjust law if they themselves were subject to that same law? This same idea can be adapted to the concept of law enforcement and one of its limbs, threat detection. If all members of society agree to act within its laws, shouldn’t we also be willing to subject to any process which upholds the law? In pursuing effective threat detection, we might build systems which step into the private lives of the people and violate their right to privacy. “The Fourth Amendment provides the primary legal framework governing AI surveillance, protecting individuals against unreasonable searches and seizures.” ~\cite{used_tenth(9)} We might build systems with deep seeded bias that end up treating certain members of society unfairly in its analysis. 
\subsection{Q-One Conclusion}
Both risks must be dealt with, but assuming they are dealt with diligently, there may yet exist a form of automation in which the pros vastly outweigh the cons. The answer to whether automation should exist in threat detection is dependent on the existence of this optimal system. If we put in the work to remove enough bias and write legislation to limit its use to already-public spheres, then yes. But if we are unwilling to make these concessions, then no.
\subsection{Question Two}
The second question is “Should we regulate the training of machine learning models, and if so, how?” Social contract theory believes in a set of rules governing how members of society should engage with one another. Regardless of the philosophical debate around machines and consciousness, the models we are building today have abilities that match and often exceed the abilities of humans within many fields. They are similar enough that training a ML model should in some respects be viewed similarly to training a human. If we would expect a human to be able to comprehend the rules of society which dictate how to treat other fairly and with respect, we should train our models to do the same.
\subsection{Q-Two Conclusion}
So if we expect our human members of society to learn and employ the standards set for all humans, we should either train our ML models with these same rules in mind or develop a specific set of rules which apply to automated systems. So the answer to this question would be “yes” and my personal take is that we should recognize the abilities and the differences of such models, and build a set of rules specifically designed for automated decision making systems.
\subsection{Question Three}
The last question we will ask in this paper is “Given what we have seen, how should we handle the existing bias in the worlds technological infrastructure?” Simply the fact that we have created and implemented automated systems that contain bias suggests a few things about the state of the world. Does this mean we have neglected to define social contracts within our societies? Does this mean we don’t consider a contract to be social if it pertains to the development of a new piece of technology? Does this mean that the world views laws governing the development and implementation of new technologies to be an infringement of some kind of negative right? For whatever reason, we are here now and this is probably due to an absence of relevant social contracts.
\subsection{Q-Three Conclusion}
In conclusion, an attempt to generate effective legislation to govern the development and implementation of automated systems that will fix existing problems and prevent future problems from arising seems unlikely. If we want to avoid a band-aid fix, and build a strong foundation for addressing these kinds of issues in the future, we will need to work together as the global society that we have been forced into to develop a framework of legislation that we can all agree on and will provide a reasonable path forward for the tech industry. This is no easy task, it will take time, and before we can even begin, we will have to set aside national differences. To answer the question, it would probably be more feasible to enact some kind of global ban on using these systems but also enact a financial relief system to ensure those affected by the sudden halt in revenue would be taken care of. After accomplishing these two things, we would hopefully have reached a relatively stable global financial market, and the harm being cause would have been stopped in its tracks. From this point, the inflammation might have subsided enough for the world’s nations to start talking about a more long-term solution.
\section{Beyond the Superficial}
\subsection{Practical Development}
As citizens, we should push for strong, clear regulation of development pipelines in all industries. It is naïve to assume our creative endeavors will yield risk free results. But we must strive to walk a fine line. Constricting creative freedom too much will stifle humanities ability to grow and change with the world. And no limitations will breed an abundance of thoughtless harm. We must seek a global balance. Technological innovation and the capacity of such creations have force humanity into a global community. We must recognize this and build a form of leadership for ourselves which can effectively consider all cultures, and write considerate, but realistic legislation to manage such a community. As a first step, we should agree to a set of tests which all new technologies must pass before they can be integrated into the new global society.
\subsection{RTIS Responsibilities}
As for myself and my degree at DigiPen, it is clear to me that I’m learning many foundational skills required of an engineer who will eventually be pushing the state of the art in software. I have specifically opted to take machine learning courses as part of my elective course load, and I now have a basic understanding of how to implement the same kinds of algorithms as are discussed in Coded Bias. Even if the world does not impose such expectations upon me, I will be seeking to diversify my perspectives on any and all projects I work on. I will seek to find the flaws or vulnerabilities in whatever systems I end up building, testing, or implementing in the world. If I can, I will try to build systems without strong bias, and without the ability to use existing bias to cause great harm.

\bibliography{poop}{}
\bibliographystyle{ACM-Reference-Format}
\end{document}